Comme indique, il est possible de garder la meme formulation en termes
de flux pour la correction de premier ordre de l'equation de Fokker-Planck
qui decrit les effets de derive. On conserve donc pour les expressions
discretes les meme notations concernant les indices. La premiere quantite
a calculer a partir de l'equation de Fokker-Planck d'ordre zero est
la fonction $\widetilde{f}^{(0)}$\[
\widetilde{f}^{(0)}(p,\xi_{0},r)=-\frac{p\xi_{0}}{q_{e}B_{\theta,0}}\frac{\partial f_{0}^{(0)}(p,\xi_{0},r)}{\partial r}\]
 qui est non-nulle en presence d'un gradient de pression. La grille
radiale pouvant etre non-uniforme, il convient de determiner precisement
la derivee radiale.

Soient $r_{-},r$ et $r_{+}$ les trois positions radiales ou sont
calculees les ditributions $f_{0}^{(0)}.$ Posons $\Delta r_{-}=r_{-}-r,$
$\Delta r_{+}=r_{+}-r$ et $\Delta r=r_{+}-r_{-}$. On considere une
interpolation parabolique de la forme $y=ar^{2}+br+c$ pour calculer
la derivee radiale qui vaut $\frac{dy}{dr}=2ar+b.$ Les coefficients
$a,b$ et $c$ etant determines par les valeurs $y_{-},y$ et $y_{+}$
aux points $r_{-},r$ et $r_{+},$ on montre alors simplement que
\[
\left.\frac{dy}{dr}\right|_{r=r_{0}}=\frac{1}{\det V}\left(-\Delta r_{+}^{2}y_{-}+\Delta r\left(\Delta r_{+}+\Delta r_{-}\right)y_{0}+\Delta r_{-}^{2}y_{+}\right)\]
 ou $\overline{\overline{V}}$ est la matrice de Van der Monde d'ordre
3 \[
\overline{\overline{V}}=\left[\begin{array}{lll}
1 & r_{-} & r_{-}^{2}\\
1 & r & r^{2}\\
1 & r_{+} & r_{+}^{2}\end{array}\right]\]
 dont le determinant vaut \[
\det V=-\Delta r_{+}\Delta r_{-}\Delta r\]


On obtient alors \[
\left.\frac{dy}{dr}\right|_{r}=\frac{\Delta r_{+}}{\Delta r\Delta r_{-}}y_{-}-\frac{\left(\Delta r_{+}+\Delta r_{-}\right)}{\Delta r_{+}\Delta r_{-}}y-\frac{\Delta r_{-}}{\Delta r\Delta r_{+}}y_{+}\]
 et pour $\widetilde{f}^{(0)},$\begin{eqnarray*}
 &  & \widetilde{f}_{0}^{(0)(k+1)}(p_{i+1/2},\xi_{0,j+1/2},r_{l+1/2})=-\frac{p\xi_{0,j+1/2}}{q_{e}B_{\theta,0,l+1/2}}\times\\
 &  & \left[\frac{\left(r_{l+3/2}-r_{l+1/2}\right)}{\left(r_{l+3/2}-r_{l-1/2}\right)\left(r_{l-1/2}-r_{l+1/2}\right)}f_{0}^{(0)(k+1)}(p_{i+1/2},\xi_{0,j+1/2},r_{l-1/2})\right.\\
 &  & -\frac{\left(r_{l+3/2}-2r_{l+1/2}+r_{l-1/2}\right)}{\left(r_{l+3/2}-r_{l+1/2}\right)\left(r_{l-1/2}-r_{l+1/2}\right)}f_{0}^{(0)(k+1)}(p_{i+1/2},\xi_{0,j+1/2},r_{l+1/2})\\
 &  & -\left.\frac{\left(r_{l-1/2}-r_{l+1/2}\right)}{\left(r_{l+3/2}-r_{l-1/2}\right)\left(r_{l+3/2}-r_{l+1/2}\right)}f_{0}^{(0)(k+1)}(p_{i+1/2},\xi_{0,j+1/2},r_{l+3/2})\right]\end{eqnarray*}


Par ailleurs, il convient de reexprimer l'equation conservative faisant
intervenir $\widetilde{f}^{(0)}$ sous forme discrete, puisque des
termes d'interpolation nouveaux vont apparaitre compte tenu des effets
de gradients. D'apres la formulation conservative universelle \[
\widetilde{\mathbf{S}}_{\mathbf{p}}^{(0)}=-\overline{\overline{\widetilde{D}}}^{(0)}\cdot\nabla_{p,\xi_{0}}\widetilde{f}^{(0)}+\overline{\widetilde{F}}^{(0)}\widetilde{f}^{(0)}\]
 les valeurs suivantes de $\widetilde{f}^{(0)}$ qui apparaissent
sur la grille entiere en $\left(p,\xi\right)$ a la position radiale
$l+1/2$ sont \[
\left\{ \begin{array}{l}
\widetilde{f}_{i+1,j+1/2}^{\left(0\right)(k+1)}=-\frac{p_{i+1}\xi_{0,j+1/2}}{q_{e}B_{\theta,0,l+1/2}}\left.\frac{\partial f_{0}^{(0)}}{\partial r}\right|_{i+1,j+1/2}^{(k+1)}\\
\widetilde{f}_{i,j+1/2}^{\left(0\right)(k+1)}=-\frac{p_{i}\xi_{0,j+1/2}}{q_{e}B_{\theta,0,l+1/2}}\left.\frac{\partial f_{0}^{(0)}}{\partial r}\right|_{i,j+1/2}^{(k+1)}\\
\widetilde{f}_{i+1/2,j+1}^{\left(0\right)(k+1)}=-\frac{p_{i+1/2}\xi_{0,j+1}}{q_{e}B_{\theta,0,l+1/2}}\left.\frac{\partial f_{0}^{(0)}}{\partial r}\right|_{i+1/2,j+1}^{(k+1)}\\
\widetilde{f}_{i+1/2,j}^{\left(0\right)(k+1)}=-\frac{p_{i+1/2}\xi_{0,j}}{q_{e}B_{\theta,0,l+1/2}}\left.\frac{\partial f_{0}^{(0)}}{\partial r}\right|_{i+1/2,j}^{(k+1)}\end{array}\right.\]


En reportant les expressions des interpolations de la distribution
$f_{0}^{(0)}$ donnees au paragraphe precedent en $l+1/2$, on obtient
donc \begin{eqnarray*}
\left.\frac{\partial f_{0}^{(0)}}{\partial r}\right|_{i+1,j+1/2}^{(k+1)} & = & \left(1-\delta_{p,i+1,j+1/2}^{(0)}\right)\left.\frac{\partial f_{0}^{(0)}}{\partial r}\right|_{i+3/2,j+1/2}^{(k+1)}+\delta_{p,i+1,j+1/2}^{(0)}\left.\frac{\partial f_{0}^{(0)}}{\partial r}\right|_{i+1/2,j+1/2}^{(k+1)}\\
 &  & -\frac{\partial\delta_{p,i+1,j+1/2}^{(0)}}{\partial r}\left[f_{0,i+3/2,j+1/2}^{\left(0\right)(k+1)}-f_{0,i+1/2,j+1/2}^{\left(0\right)(k+1)}\right]\end{eqnarray*}
\begin{eqnarray*}
\left.\frac{\partial f_{0}^{(0)}}{\partial r}\right|_{i,j+1/2}^{(k+1)} & = & \left(1-\delta_{p,i,j+1/2}^{(0)}\right)\left.\frac{\partial f_{0}^{(0)}}{\partial r}\right|_{i+1/2,j+1/2}^{(k+1)}+\delta_{p,i,j+1/2}^{(0)}\left.\frac{\partial f_{0}^{(0)}}{\partial r}\right|_{i-1/2,j+1/2}^{(k+1)}\\
 &  & -\frac{\partial\delta_{p,i,j+1/2}^{(0)}}{\partial r}\left[f_{0,i+1/2,j+1/2}^{\left(0\right)(k+1)}-f_{0,i-1/2,j+1/2}^{\left(0\right)(k+1)}\right]\end{eqnarray*}
\begin{eqnarray*}
\left.\frac{\partial f_{0}^{(0)}}{\partial r}\right|_{i+1/2,j+1}^{(k+1)} & = & \left(1-\delta_{\xi,i+1/2,j+1}^{(0)}\right)\left.\frac{\partial f_{0}^{(0)}}{\partial r}\right|_{i+1/2,j+3/2}^{(k+1)}+\delta_{\xi,i+1/2,j+1}^{(0)}\left.\frac{\partial f_{0}^{(0)}}{\partial r}\right|_{i+1/2,j+1/2}^{(k+1)}\\
 &  & -\frac{\partial\delta_{\xi,i+1/2,j+1}^{(0)}}{\partial r}\left[f_{0,i+1/2,j+3/2}^{\left(0\right)(k+1)}-f_{0,i+1/2,j+1/2}^{\left(0\right)(k+1)}\right]\end{eqnarray*}
\begin{eqnarray*}
\left.\frac{\partial f_{0}^{(0)}}{\partial r}\right|_{i+1/2,j}^{(k+1)} & = & \left(1-\delta_{\xi,i+1/2,j}^{(0)}\right)\left.\frac{\partial f_{0}^{(0)}}{\partial r}\right|_{i+1/2,j+1/2}^{(k+1)}+\delta_{\xi,i+1/2,j}^{(0)}\left.\frac{\partial f_{0}^{(0)}}{\partial r}\right|_{i+1/2,j-1/2}^{(k+1)}\\
 &  & -\frac{\partial\delta_{\xi,i+1/2,j}^{(0)}}{\partial r}\left[f_{0,i+1/2,j+1/2}^{\left(0\right)(k+1)}-f_{0,i+1/2,j-1/2}^{\left(0\right)(k+1)}\right]\end{eqnarray*}


En multipliant les derivees par le coefficient $-\frac{p\xi_{0}}{q_{e}B_{\theta,0}}$
a la position discrete appropriee, on trouve au point de grille $\left(i+1,j+1/2\right),$\begin{eqnarray*}
\widetilde{f}_{i+1,j+1/2}^{\left(0\right)(k+1)} & = & -\frac{p_{i+1}\xi_{0,j+1/2}}{q_{e}B_{\theta,0,l+1/2}}\times\left[\left(1-\delta_{p,i+1,j+1/2}^{(0)}\right)\left.\frac{\partial f_{0}^{(0)}}{\partial r}\right|_{i+3/2,j+1/2}^{(k+1)}\right.\\
 &  & +\delta_{p,i+1,j+1/2}^{(0)}\left.\frac{\partial f_{0}^{(0)}}{\partial r}\right|_{i+1/2,j+1/2}^{(k+1)}\\
 &  & -\frac{\partial\delta_{p,i+1,j+1/2}^{(0)}}{\partial r}f_{0,i+3/2,j+1/2}^{\left(0\right)(k+1)}\\
 &  & \left.+\frac{\partial\delta_{p,i+1,j+1/2}^{(0)}}{\partial r}f_{0,i+1/2,j+1/2}^{\left(0\right)(k+1)}\right]\end{eqnarray*}
 soit \begin{eqnarray*}
\widetilde{f}_{i+1,j+1/2}^{\left(0\right)(k+1)} & = & \left(1-\delta_{p,i+1,j+1/2}^{(0)}\right)\frac{p_{i+1}}{p_{i+3/2}}\widetilde{f}_{i+3/2,j+1/2}^{\left(0\right)(k+1)}\\
 &  & +\delta_{p,i+1,j+1/2}^{(0)}\frac{p_{i+1}}{p_{i+1/2}}\widetilde{f}_{i+1/2,j+1/2}^{\left(0\right)(k+1)}\\
 &  & +\frac{p_{i+1}\xi_{0,j+1/2}}{q_{e}B_{\theta,0,l+1/2}}\frac{\partial\delta_{p,i+1,j+1/2}^{(0)}}{\partial r}f_{0,i+3/2,j+1/2}^{\left(0\right)(k+1)}\\
 &  & -\frac{p_{i+1}\xi_{0,j+1/2}}{q_{e}B_{\theta,0,l+1/2}}\frac{\partial\delta_{p,i+1,j+1/2}^{(0)}}{\partial r}f_{0,i+1/2,j+1/2}^{\left(0\right)(k+1)}\end{eqnarray*}
 , au point de grille $\left(i,j+1/2\right),$\begin{eqnarray*}
\widetilde{f}_{i,j+1/2}^{\left(0\right)(k+1)} & = & -\frac{p_{i}\xi_{0,j+1/2}}{q_{e}B_{\theta,0,l+1/2}}\times\left[\left(1-\delta_{p,i,j+1/2}^{(0)}\right)\left.\frac{\partial f_{0}^{(0)}}{\partial r}\right|_{i+1/2,j+1/2}^{(k+1)}\right.\\
 &  & +\delta_{p,i,j+1/2}^{(0)}\left.\frac{\partial f_{0}^{(0)}}{\partial r}\right|_{i-1/2,j+1/2}^{(k+1)}\\
 &  & -\frac{\partial\delta_{p,i,j+1/2}^{(0)}}{\partial r}f_{0,i+1/2,j+1/2}^{\left(0\right)(k+1)}\\
 &  & \left.+\frac{\partial\delta_{p,i,j+1/2}^{(0)}}{\partial r}f_{0,i-1/2,j+1/2}^{\left(0\right)(k+1)}\right]\end{eqnarray*}
 soit \begin{eqnarray*}
\widetilde{f}_{i,j+1/2}^{\left(0\right)(k+1)} & = & \left(1-\delta_{p,i,j+1/2}^{(0)}\right)\frac{p_{i}}{p_{i+1/2}}\widetilde{f}_{i+1/2,j+1/2}^{\left(0\right)(k+1)}\\
 &  & +\delta_{p,i,j+1/2}^{(0)}\frac{p_{i}}{p_{i-1/2}}\widetilde{f}_{i-1/2,j+1/2}^{\left(0\right)(k+1)}\\
 &  & +\frac{p_{i}\xi_{0,j+1/2}}{q_{e}B_{\theta,0,l+1/2}}\frac{\partial\delta_{p,i,j+1/2}^{(0)}}{\partial r}f_{0,i+1/2,j+1/2}^{\left(0\right)(k+1)}\\
 &  & -\frac{p_{i}\xi_{0,j+1/2}}{q_{e}B_{\theta,0,l+1/2}}\frac{\partial\delta_{p,i,j+1/2}^{(0)}}{\partial r}f_{0,i-1/2,j+1/2}^{\left(0\right)(k+1)}\end{eqnarray*}
 , au point de grille $\left(i+1/2,j+1\right),$\begin{eqnarray*}
\widetilde{f}_{i+1/2,j+1}^{\left(0\right)(k+1)} & = & -\frac{p_{i+1/2}\xi_{0,j+1}}{q_{e}B_{\theta,0,l+1/2}}\times\left[\left(1-\delta_{\xi,i+1/2,j+1}^{(0)}\right)\left.\frac{\partial f_{0}^{(0)}}{\partial r}\right|_{i+1/2,j+3/2}^{(k+1)}\right.\\
 &  & +\delta_{\xi,i+1/2,j+1}^{(0)}\left.\frac{\partial f_{0}^{(0)}}{\partial r}\right|_{i+1/2,j+1/2}^{(k+1)}\\
 &  & -\frac{\partial\delta_{\xi,i+1/2,j+1}^{(0)}}{\partial r}f_{0,i+1/2,j+3/2}^{\left(0\right)(k+1)}\\
 &  & \left.+\frac{\partial\delta_{\xi,i+1/2,j+1}^{(0)}}{\partial r}f_{0,i+1/2,j+1/2}^{\left(0\right)(k+1)}\right]\end{eqnarray*}
 soit \begin{eqnarray*}
\widetilde{f}_{i+1/2,j+1}^{\left(0\right)(k+1)} & = & \left(1-\delta_{\xi,i+1/2,j+1}^{(0)}\right)\frac{\xi_{0,j+1}}{\xi_{0,j+3/2}}\widetilde{f}_{i+1/2,j+3/2}^{\left(0\right)(k+1)}\\
 &  & +\delta_{\xi,i+1/2,j+1}^{(0)}\frac{\xi_{0,j+1}}{\xi_{0,j+1/2}}\widetilde{f}_{i+1/2,j+1/2}^{\left(0\right)(k+1)}\\
 &  & +\frac{p_{i+1/2}\xi_{0,j+1}}{q_{e}B_{\theta,0,l+1/2}}\frac{\partial\delta_{\xi,i+1/2,j+1}^{(0)}}{\partial r}f_{0,i+1/2,j+3/2}^{\left(0\right)(k+1)}\\
 &  & -\frac{p_{i+1/2}\xi_{0,j+1}}{q_{e}B_{\theta,0,l+1/2}}\frac{\partial\delta_{\xi,i+1/2,j+1}^{(0)}}{\partial r}f_{0,i+1/2,j+1/2}^{\left(0\right)(k+1)}\end{eqnarray*}
 et finalement en $\left(i+1/2,j\right),$\begin{eqnarray*}
\widetilde{f}_{i+1/2,j}^{\left(0\right)(k+1)} & = & -\frac{p_{i+1/2}\xi_{0,j}}{q_{e}B_{\theta,0,l+1/2}}\times\left[\left(1-\delta_{\xi,i+1/2,j}^{(0)}\right)\left.\frac{\partial f_{0}^{(0)}}{\partial r}\right|_{i+1/2,j+1/2}^{(k+1)}\right.\\
 &  & +\delta_{\xi,i+1/2,j}^{(0)}\left.\frac{\partial f_{0}^{(0)}}{\partial r}\right|_{i+1/2,j-1/2}^{(k+1)}\\
 &  & -\frac{\partial\delta_{\xi,i+1/2,j}^{(0)}}{\partial r}f_{0,i+1/2,j+1/2}^{\left(0\right)(k+1)}\\
 &  & \left.+\frac{\partial\delta_{\xi,i+1/2,j}^{(0)}}{\partial r}f_{0,i+1/2,j-1/2}^{\left(0\right)(k+1)}\right]\end{eqnarray*}
 soit \begin{eqnarray*}
\widetilde{f}_{i+1/2,j}^{\left(0\right)(k+1)} & = & \left(1-\delta_{\xi,i+1/2,j}^{(0)}\right)\frac{\xi_{0,j}}{\xi_{0,j+3/2}}\widetilde{f}_{i+1/2,j+1/2}^{\left(0\right)(k+1)}\\
 &  & +\delta_{\xi,i+1/2,j}^{(0)}\frac{\xi_{0,j}}{\xi_{0,j+1/2}}\widetilde{f}_{i+1/2,j-1/2}^{\left(0\right)(k+1)}\\
 &  & +\frac{p_{i+1/2}\xi_{0,j}}{q_{e}B_{\theta,0,l+1/2}}\frac{\partial\delta_{\xi,i+1/2,j}^{(0)}}{\partial r}f_{0,i+1/2,j+1/2}^{\left(0\right)(k+1)}\\
 &  & -\frac{p_{i+1/2}\xi_{0,j}}{q_{e}B_{\theta,0,l+1/2}}\frac{\partial\delta_{\xi,i+1/2,j}^{(0)}}{\partial r}f_{0,i+1/2,j-1/2}^{\left(0\right)(k+1)}\end{eqnarray*}


Les derivees radiales des coefficients de Chang et Cooper sont calcules
de la meme maniere que pour celles de la fonction de distribution
$f_{0}^{\left(0\right)},$ ce qui donne \begin{eqnarray*}
\left.\frac{\partial\delta_{p,i+1,j+1/2}^{(0)}}{\partial r}\right|_{l+1/2} & = & \left[\frac{\left(r_{l+3/2}-r_{l+1/2}\right)}{\left(r_{l+3/2}-r_{l-1/2}\right)\left(r_{l-1/2}-r_{l+1/2}\right)}\delta_{p,i+1,j+1/2}^{(0)}\left(r_{l-1/2}\right)\right.\\
 &  & -\frac{\left(r_{l+3/2}-2r_{l+1/2}+r_{l-1/2}\right)}{\left(r_{l+3/2}-r_{l+1/2}\right)\left(r_{l-1/2}-r_{l+1/2}\right)}\delta_{p,i+1,j+1/2}^{(0)}\left(r_{l+1/2}\right)\\
 &  & -\left.\frac{\left(r_{l-1/2}-r_{l+1/2}\right)}{\left(r_{l+3/2}-r_{l-1/2}\right)\left(r_{l+3/2}-r_{l+1/2}\right)}\delta_{p,i+1,j+1/2}^{(0)}\left(r_{l+3/2}\right)\right]\end{eqnarray*}
\begin{eqnarray*}
\left.\frac{\partial\delta_{p,i,j+1/2}^{(0)}}{\partial r}\right|_{l+1/2} & = & \left[\frac{\left(r_{l+3/2}-r_{l+1/2}\right)}{\left(r_{l+3/2}-r_{l-1/2}\right)\left(r_{l-1/2}-r_{l+1/2}\right)}\delta_{p,i,j+1/2}^{(0)}\left(r_{l-1/2}\right)\right.\\
 &  & -\frac{\left(r_{l+3/2}-2r_{l+1/2}+r_{l-1/2}\right)}{\left(r_{l+3/2}-r_{l+1/2}\right)\left(r_{l-1/2}-r_{l+1/2}\right)}\delta_{p,i,j+1/2}^{(0)}\left(r_{l+1/2}\right)\\
 &  & -\left.\frac{\left(r_{l-1/2}-r_{l+1/2}\right)}{\left(r_{l+3/2}-r_{l-1/2}\right)\left(r_{l+3/2}-r_{l+1/2}\right)}\delta_{p,i,j+1/2}^{(0)}\left(r_{l+3/2}\right)\right]\end{eqnarray*}
\begin{eqnarray*}
\left.\frac{\partial\delta_{\xi,i+1/2,j+1}^{(0)}}{\partial r}\right|_{l+1/2} & = & \left[\frac{\left(r_{l+3/2}-r_{l+1/2}\right)}{\left(r_{l+3/2}-r_{l-1/2}\right)\left(r_{l-1/2}-r_{l+1/2}\right)}\delta_{\xi,i+1/2,j+1}^{(0)}\left(r_{l-1/2}\right)\right.\\
 &  & -\frac{\left(r_{l+3/2}-2r_{l+1/2}+r_{l-1/2}\right)}{\left(r_{l+3/2}-r_{l+1/2}\right)\left(r_{l-1/2}-r_{l+1/2}\right)}\delta_{\xi,i+1/2,j+1}^{(0)}\left(r_{l+1/2}\right)\\
 &  & -\left.\frac{\left(r_{l-1/2}-r_{l+1/2}\right)}{\left(r_{l+3/2}-r_{l-1/2}\right)\left(r_{l+3/2}-r_{l+1/2}\right)}\delta_{\xi,i+1/2,j+1}^{(0)}\left(r_{l+3/2}\right)\right]\end{eqnarray*}
\begin{eqnarray*}
\left.\frac{\partial\delta_{\xi,i+1/2,j}^{(0)}}{\partial r}\right|_{l+1/2} & = & \left[\frac{\left(r_{l+3/2}-r_{l+1/2}\right)}{\left(r_{l+3/2}-r_{l-1/2}\right)\left(r_{l-1/2}-r_{l+1/2}\right)}\delta_{\xi,i+1/2,j}^{(0)}\left(r_{l-1/2}\right)\right.\\
 &  & -\frac{\left(r_{l+3/2}-2r_{l+1/2}+r_{l-1/2}\right)}{\left(r_{l+3/2}-r_{l+1/2}\right)\left(r_{l-1/2}-r_{l+1/2}\right)}\delta_{\xi,i+1/2,j}^{(0)}\left(r_{l+1/2}\right)\\
 &  & -\left.\frac{\left(r_{l-1/2}-r_{l+1/2}\right)}{\left(r_{l+3/2}-r_{l-1/2}\right)\left(r_{l+3/2}-r_{l+1/2}\right)}\delta_{\xi,i+1/2,j}^{(0)}\left(r_{l+3/2}\right)\right]\end{eqnarray*}
